\section{Week 8: Midterm Review: The Whole Pipeline in One Case Study}

\begin{quote}
\textbf{Chapters 1--3 review} · Lecture 1 h (integrated case study) ·
Lab slot: \textbf{midterm exam} (90 min) + project consultation (30 min)
\end{quote}

\begin{center}\rule{0.5\linewidth}{0.5pt}\end{center}

\subsection{Key terms}

Important terms of this lesson, with their Persian equivalents (repeated
even if they appeared in earlier weeks).

\begin{longtable}[]{@{}
  >{\raggedright\arraybackslash}p{(\linewidth - 4\tabcolsep) * \real{0.3333}}
  >{\raggedright\arraybackslash}p{(\linewidth - 4\tabcolsep) * \real{0.3333}}
  >{\raggedright\arraybackslash}p{(\linewidth - 4\tabcolsep) * \real{0.3333}}@{}}
\caption{\centering Key terms}\tabularnewline
\toprule\noalign{}
\begin{minipage}[b]{\linewidth}\raggedright
English
\end{minipage} & \begin{minipage}[b]{\linewidth}\raggedright
\RL{\CMfa فارسی}
\end{minipage} & \begin{minipage}[b]{\linewidth}\raggedright
\RL{\CMfa توضیح کوتاه}
\end{minipage} \\
\midrule\noalign{}
\endfirsthead
\toprule\noalign{}
\begin{minipage}[b]{\linewidth}\raggedright
English
\end{minipage} & \begin{minipage}[b]{\linewidth}\raggedright
\RL{\CMfa فارسی}
\end{minipage} & \begin{minipage}[b]{\linewidth}\raggedright
\RL{\CMfa توضیح کوتاه}
\end{minipage} \\
\midrule\noalign{}
\endhead
\bottomrule\noalign{}
\endlastfoot
formal model / formalize & \RL{\CMfa مدل فرمال} / \RL{\CMfa فرمال کردن} & \RL{\CMfa نوشتن
نظریه به شکل معادله یا الگوریتم} \\
simulation & \RL{\CMfa شبیه‌سازی} & \RL{\CMfa تولید داده‌ی مصنوعی با اجرای مدل} \\
fit / fitting & \RL{\CMfa برازش} & \RL{\CMfa تنظیم پارامترها تا مدل با داده جور
شود} \\
parameter recovery & \RL{\CMfa بازیابی پارامتر} & \RL{\CMfa آیا برازش، پارامترهای
معلوم شبیه‌سازی را برمی‌گرداند} \\
likelihood & \RL{\CMfa درست‌نمایی} & \RL{\CMfa احتمال داده به‌عنوان تابعی از
پارامترها} \\
criterion (c) & \RL{\CMfa ملاک} & \RL{\CMfa آستانه‌ی تصمیم؛ نشان‌دهنده‌ی سوگیری
پاسخ} \\
least squares & \RL{\CMfa کمترین مربعات} & \RL{\CMfa روش برازش با کمینه کردن مجموع
مربع خطاها} \\
cross-validation (CV) & \RL{\CMfa اعتبارسنجی متقابل} (cross-validation) &
\RL{\CMfa برازش روی بخشی از داده و آزمون روی بخش کنارگذاشته} \\
sensitivity (d′) & \RL{\CMfa حساسیت} & \RL{\CMfa فاصله‌ی توزیع علامت و نویز} \\
parameter & \RL{\CMfa پارامتر} & \RL{\CMfa کمیت قابل تنظیم مدل که بین افراد یا شرایط
فرق می‌کند} \\
SSE / RMSD & \RL{\CMfa مجموع مربعات خطا} / \RL{\CMfa ریشه‌ی میانگین مربع انحراف} &
\RL{\CMfa معیارهای فاصله‌ی داده و پیش‌بینی} \\
log-likelihood / NLL & \RL{\CMfa لگاریتم درست‌نمایی} / \RL{\CMfa منفی آن} &
\RL{\CMfa مقداری که بهینه‌سازها کمینه یا بیشینه می‌کنند} \\
estimate / estimator & \RL{\CMfa برآورد} / \RL{\CMfa برآوردگر} & \RL{\CMfa مقدار به‌دست‌آمده
برای پارامتر} / \RL{\CMfa روش به دست آوردن آن} \\
noise & \RL{\CMfa نویز} & \RL{\CMfa تغییرات تصادفی در شواهد یا در داده} \\
deviance (G²) & \RL{\CMfa انحراف} (deviance) & \RL{\CMfa معیار عدم برازش نسبت به مدل
اشباع‌شده} \\
AIC / BIC & \RL{\CMfa معیار اطلاعاتی آکائیکه} / \RL{\CMfa بیزی} & \RL{\CMfa برازش به‌علاوه‌ی
جریمه برای تعداد پارامتر} \\
research design & \RL{\CMfa طرح آزمایش} & \RL{\CMfa شرایط، تعداد کوشش‌ها و
نمونه‌ها} \\
confusion matrix & \RL{\CMfa ماتریس درهم‌ریختگی} (confusion matrix) & P(\RL{\CMfa مدل
انتخاب‌شده} \\
hit & \RL{\CMfa اصابت} (hit) & \RL{\CMfa پاسخ «بله» وقتی علامت وجود دارد} \\
response time (RT) & \RL{\CMfa زمان واکنش} & \RL{\CMfa فاصله‌ی ارائه‌ی محرک تا
پاسخ} \\
signal detection theory (SDT) & \RL{\CMfa نظریه‌ی تشخیص علامت} & \RL{\CMfa مدلی که
حساسیت را از سوگیری پاسخ جدا می‌کند} \\
evidence & \RL{\CMfa شواهد} & \RL{\CMfa اطلاعاتی که برای تصمیم جمع می‌شود} \\
exponential / power function & \RL{\CMfa تابع نمایی} / \RL{\CMfa توانی} & \RL{\CMfa دو شکل
رایج برای منحنی فراموشی و یادگیری} \\
training / test data & \RL{\CMfa داده‌ی آموزش} / \RL{\CMfa آزمون} & \RL{\CMfa داده‌ی برازش}
/ \RL{\CMfa داده‌ی کنارگذاشته برای سنجش} \\
\end{longtable}

\begin{center}\rule{0.5\linewidth}{0.5pt}\end{center}

\subsection{Learning outcomes}

By the end of this week, students will be able to:

\begin{enumerate}
\def\labelenumi{\arabic{enumi}.}
\tightlist
\item
  \textbf{Execute} the complete modeling pipeline (formalize \ensuremath{\rightarrow} simulate
  \ensuremath{\rightarrow} fit \ensuremath{\rightarrow} evaluate \ensuremath{\rightarrow} recover \ensuremath{\rightarrow} interpret) on a new problem without
  step-by-step guidance.
\item
  \textbf{Select} an appropriate likelihood, fitting method and
  evaluation method for a given dataset, and \textbf{justify} each
  choice.
\item
  \textbf{Interpret} a set of competing model fits, information criteria
  and recovery results, and \textbf{write} a defensible conclusion that
  states its limits.
\item
  \textbf{Recall} and \textbf{apply} the key formulas of Weeks 1--7 (the
  midterm formula sheet).
\end{enumerate}

\begin{center}\rule{0.5\linewidth}{0.5pt}\end{center}

\subsection{Motivation \& big picture}

Over seven weeks we built a toolkit piece by piece: formal models (Weeks
1--2), least squares and optimization (Week 3), maximum likelihood (Week
4), uncertainty and multiple participants (Week 5), cross-validation
(Week 6), information criteria and recovery (Week 7). Each lab practiced
one or two pieces. Real research needs all of them at once, in the right
order, with judgment about what matters.

Today's lecture runs the entire pipeline on a \textbf{new} problem: does
studying a word more often make it more \emph{memorable} (higher
sensitivity), and if so, how does sensitivity grow with repetitions? The
case study is also the best preparation for the midterm, which asks you
to make and justify the same kinds of decisions.

\begin{center}\rule{0.5\linewidth}{0.5pt}\end{center}

\Needspace{11\baselineskip}

\subsection{Lecture notes}

\paragraph{In plain words}

This week puts all the pieces of the first half together on one new
question: \emph{does studying a word more often make it easier to
recognize later?} We write the theory as an SDT model in which
sensitivity \(d'\) depends on the number of repetitions. We propose
three versions of that dependence (no effect, a smooth increase, any
pattern), fit all three by maximum likelihood, check whether each
misfits (\(G^2\)), compare them (LRT, AIC, BIC), check that the method
works at this design (recovery), and finally interpret the winner in
plain psychological language, including what the data \emph{cannot} tell
us. If you can follow this one example from start to end, you understand
the first half of the course.

\textbf{Suggested timing (60 min):} concept map 10 · case study 40 ·
exam briefing 10.

\subsubsection{1. Concept map of Weeks 1--7}

\begin{longtable}[]{@{}
  >{\raggedright\arraybackslash}p{(\linewidth - 6\tabcolsep) * \real{0.2500}}
  >{\raggedright\arraybackslash}p{(\linewidth - 6\tabcolsep) * \real{0.2500}}
  >{\raggedright\arraybackslash}p{(\linewidth - 6\tabcolsep) * \real{0.2500}}
  >{\raggedright\arraybackslash}p{(\linewidth - 6\tabcolsep) * \real{0.2500}}@{}}
\caption{\centering 1. Concept map of Weeks 1--7}\tabularnewline
\toprule\noalign{}
\begin{minipage}[b]{\linewidth}\raggedright
Pipeline step
\end{minipage} & \begin{minipage}[b]{\linewidth}\raggedright
Question
\end{minipage} & \begin{minipage}[b]{\linewidth}\raggedright
Tools
\end{minipage} & \begin{minipage}[b]{\linewidth}\raggedright
Week
\end{minipage} \\
\midrule\noalign{}
\endfirsthead
\toprule\noalign{}
\begin{minipage}[b]{\linewidth}\raggedright
Pipeline step
\end{minipage} & \begin{minipage}[b]{\linewidth}\raggedright
Question
\end{minipage} & \begin{minipage}[b]{\linewidth}\raggedright
Tools
\end{minipage} & \begin{minipage}[b]{\linewidth}\raggedright
Week
\end{minipage} \\
\midrule\noalign{}
\endhead
\bottomrule\noalign{}
\endlastfoot
Formalize & What exactly does the theory claim? & Five components of a
model; parameters as psychological constructs & 1--2 \\
Simulate & What can (and can't) the model do? & Forward simulation,
analytic checks, prediction plots & 1--2 \\
Fit & Which parameter values best account for the data? & SSE / least
squares; NLL / MLE; grid, Nelder--Mead, L-BFGS-B; multi-start; bounds
and reparameterization & 3--4 \\
Quantify uncertainty & How precise are the estimates? & Hessian (Wald)
SEs, profile likelihood, bootstrap & 5 \\
Evaluate (absolute) & Does the model misfit beyond noise? & Residuals,
\(G^2\) & 2--4 \\
Evaluate (relative) & Which model generalizes best? & CV, LRT,
AIC/AICc/BIC, weights & 6--7 \\
Recover & Can the method work at this design? & Parameter recovery,
model recovery (confusion and inversion matrices) & 7 \\
Interpret & What does it mean psychologically, and what can't we
conclude? & Parameter meaning, design limits, aggregation artifacts &
all \\
\end{longtable}

\textbf{Decision guide for the likelihood:}

\begin{longtable}[]{@{}
  >{\raggedright\arraybackslash}p{(\linewidth - 4\tabcolsep) * \real{0.3333}}
  >{\raggedright\arraybackslash}p{(\linewidth - 4\tabcolsep) * \real{0.3333}}
  >{\raggedright\arraybackslash}p{(\linewidth - 4\tabcolsep) * \real{0.3333}}@{}}
\caption{\centering 1. Concept map of Weeks 1--7}\tabularnewline
\toprule\noalign{}
\begin{minipage}[b]{\linewidth}\raggedright
Data
\end{minipage} & \begin{minipage}[b]{\linewidth}\raggedright
Natural noise model
\end{minipage} & \begin{minipage}[b]{\linewidth}\raggedright
Week
\end{minipage} \\
\midrule\noalign{}
\endfirsthead
\toprule\noalign{}
\begin{minipage}[b]{\linewidth}\raggedright
Data
\end{minipage} & \begin{minipage}[b]{\linewidth}\raggedright
Natural noise model
\end{minipage} & \begin{minipage}[b]{\linewidth}\raggedright
Week
\end{minipage} \\
\midrule\noalign{}
\endhead
\bottomrule\noalign{}
\endlastfoot
Counts out of \(n\) (recall, hits, false alarms) & Binomial & 2, 4 \\
Single binary choices & Bernoulli & 4 \\
Continuous, roughly symmetric, equal variance (e.g., block-mean RTs) &
Gaussian (= least squares) & 4--5 \\
\end{longtable}

\subsubsection{2. Case study: does repetition increase memory
sensitivity?}

\textbf{Design.} In a recognition-memory experiment, words are studied
1, 2 or 4 times. At test, 100 old words per repetition condition and 100
new words (lures) are presented, and the participant says ``old'' or
``new''. Because lures were never studied, there is \textbf{one}
false-alarm rate, shared by all conditions, and \textbf{three} hit
rates.

\textbf{Formalization (SDT with an absolute criterion).} New words
produce evidence \(X \sim \mathcal{N}(0, 1)\); old words studied \(r\)
times produce \(X \sim \mathcal{N}(d'_r, 1)\); the participant says
``old'' if \(X > \kappa\), with one criterion \(\kappa\) for all test
items (they are intermixed, so the participant cannot know an item's
condition). Then

\[
H_r = \Phi(d'_r - \kappa), \qquad F = \Phi(-\kappa).
\]

\textbf{Three candidate models of how \(d'\) depends on repetitions:}

\begin{longtable}[]{@{}
  >{\raggedright\arraybackslash}p{(\linewidth - 6\tabcolsep) * \real{0.2500}}
  >{\raggedright\arraybackslash}p{(\linewidth - 6\tabcolsep) * \real{0.2500}}
  >{\raggedright\arraybackslash}p{(\linewidth - 6\tabcolsep) * \real{0.2500}}
  >{\raggedright\arraybackslash}p{(\linewidth - 6\tabcolsep) * \real{0.2500}}@{}}
\caption{\centering 2. Case study: does repetition increase memory sensitivity?}\tabularnewline
\toprule\noalign{}
\begin{minipage}[b]{\linewidth}\raggedright
Model
\end{minipage} & \begin{minipage}[b]{\linewidth}\raggedright
\(d'_r\)
\end{minipage} & \begin{minipage}[b]{\linewidth}\raggedright
\(k\)
\end{minipage} & \begin{minipage}[b]{\linewidth}\raggedright
Claim
\end{minipage} \\
\midrule\noalign{}
\endfirsthead
\toprule\noalign{}
\begin{minipage}[b]{\linewidth}\raggedright
Model
\end{minipage} & \begin{minipage}[b]{\linewidth}\raggedright
\(d'_r\)
\end{minipage} & \begin{minipage}[b]{\linewidth}\raggedright
\(k\)
\end{minipage} & \begin{minipage}[b]{\linewidth}\raggedright
Claim
\end{minipage} \\
\midrule\noalign{}
\endhead
\bottomrule\noalign{}
\endlastfoot
\(\mathcal{M}_0\) (no effect) & \(\delta_0\) & 2 & Repetition does not
change sensitivity \\
\(\mathcal{M}_1\) (log-linear) & \(\delta_0 + \delta_1 \log_2 r\) & 3 &
Each \emph{doubling} of study adds the same sensitivity \\
\(\mathcal{M}_2\) (free) & \(d'_1, d'_2, d'_4\) & 4 & Any pattern
(saturated: 4 parameters for 4 rates) \\
\end{longtable}

\(\mathcal{M}_0 \subset \mathcal{M}_1 \subset \mathcal{M}_2\), so all
pairs are nested.

\textbf{Likelihood.} Binomial for each of the four counts (three hit
counts, one false-alarm count).

Students run it at home as exam preparation; the lecture walks through
the outputs.

\paragraph{Reading the math}

\begin{itemize}
\tightlist
\item
  \(H_r = \Phi(d'_r - \kappa)\) and \(F = \Phi(-\kappa)\): the same SDT
  as Week 2, but the criterion is written as an \textbf{absolute}
  position \(\kappa\) on the evidence axis (measured from the center of
  the ``new word'' curve) instead of relative to the midpoint. New
  words: curve at 0, so \(F\) is the area right of \(\kappa\),
  i.e.~\(\Phi(-\kappa)\). Old words studied \(r\) times: curve at
  \(d'_r\), so \(H_r = \Phi(d'_r - \kappa)\).
\item
  One \(\kappa\) for all conditions: all test words are mixed together,
  so the participant cannot use a different threshold for words studied
  once and words studied four times.
\item
  \(\log_2 r\): the base-2 logarithm, so \(\log_2 1 = 0\),
  \(\log_2 2 = 1\), \(\log_2 4 = 2\). In model \(\mathcal{M}_1\),
  \(d'_r = \delta_0 + \delta_1\log_2 r\) means: \(\delta_0\) is the
  sensitivity after one study, and every \emph{doubling} of study adds
  \(\delta_1\).
\item
  \textbf{Numbers from the worked example:} \(\hat\kappa = 1.04\) gives
  \(F = \Phi(-1.04) = 0.15\) \ensuremath{\checkmark}. With \(\hat d'_1 = 0.82\):
  \(H_1 = \Phi(0.82 - 1.04) = \Phi(-0.22) = 0.41\), close to the
  observed 0.40 \ensuremath{\checkmark}.
\end{itemize}

\paragraph{For an engineer}

The pipeline of the concept map is the \textbf{system identification}
workflow you know: choose candidate model structures, estimate
parameters (ML), validate (residual tests \ensuremath{\approx} \(G^2\)), select the model
order (LRT/AIC/BIC \ensuremath{\approx} GLRT/AIC/MDL), and verify the procedure on
simulated data with known truth (recovery). Two differences deserve
emphasis in class: the ``system'' is a person whose parameters are
themselves the scientific result, and the amount of data is tiny by
engineering standards (400 binary responses), so uncertainty and
recovery matter much more.

\paragraph{For the instructor}

\begin{itemize}
\tightlist
\item
  \textbf{Use the concept-map table as the backbone of the session.} For
  each row, point to the corresponding output of the case study
  (e.g.~``Evaluate (absolute)'' \ensuremath{\rightarrow} the \(G^2\) column). Students see that
  every tool has a place.
\item
  \textbf{Ask the class to interpret before you do.} Show the results
  table without the interpretation paragraph and let small groups write
  two sentences: ``what do we conclude, and what can't we conclude?''
  Then compare with the text. This is exactly the skill the midterm
  tests.
\item
  \textbf{Exam briefing:} tell students that the exam rewards
  \emph{justified decisions} (which likelihood, which criterion, why
  recovery matters) more than computation. Show one past-style question
  from the ``Check your understanding'' section.
\item
  \textbf{Project consultation (lab slot):} ask each team three
  questions: What is your cognitive question? Which parameter or model
  comparison answers it? Can your design recover it? Most weak proposals
  fail on the third.
\end{itemize}

\begin{center}\rule{0.5\linewidth}{0.5pt}\end{center}

\subsection{Worked example}

\(\delta_0 = 0.8\) · \(\delta_1 = 0.6\) · \(\kappa = 1.0\) Numbers are
rounded.

\textbf{Step 1: data.}

\begin{longtable}[]{@{}llll@{}}
\caption{\centering Worked example}\tabularnewline
\toprule\noalign{}
Condition & Old words & ``Old'' responses & Rate \\
\midrule\noalign{}
\endfirsthead
\toprule\noalign{}
Condition & Old words & ``Old'' responses & Rate \\
\midrule\noalign{}
\endhead
\bottomrule\noalign{}
\endlastfoot
\(r = 1\) & 100 & 40 & 0.40 \\
\(r = 2\) & 100 & 66 & 0.66 \\
\(r = 4\) & 100 & 80 & 0.80 \\
New (lures) & 100 & 15 & 0.15 \\
\end{longtable}

\textbf{Step 2: fit and evaluate} (MLE, \(N = 400\) responses):

\begin{landscape}\small
\begin{longtable}[]{@{}
  >{\raggedright\arraybackslash}p{(\linewidth - 18\tabcolsep) * \real{0.1000}}
  >{\raggedright\arraybackslash}p{(\linewidth - 18\tabcolsep) * \real{0.1000}}
  >{\raggedright\arraybackslash}p{(\linewidth - 18\tabcolsep) * \real{0.1000}}
  >{\raggedright\arraybackslash}p{(\linewidth - 18\tabcolsep) * \real{0.1000}}
  >{\raggedright\arraybackslash}p{(\linewidth - 18\tabcolsep) * \real{0.1000}}
  >{\raggedright\arraybackslash}p{(\linewidth - 18\tabcolsep) * \real{0.1000}}
  >{\raggedright\arraybackslash}p{(\linewidth - 18\tabcolsep) * \real{0.1000}}
  >{\raggedright\arraybackslash}p{(\linewidth - 18\tabcolsep) * \real{0.1000}}
  >{\raggedright\arraybackslash}p{(\linewidth - 18\tabcolsep) * \real{0.1000}}
  >{\raggedright\arraybackslash}p{(\linewidth - 18\tabcolsep) * \real{0.1000}}@{}}
\caption{\centering Worked example}\tabularnewline
\toprule\noalign{}
\begin{minipage}[b]{\linewidth}\raggedright
Model
\end{minipage} & \begin{minipage}[b]{\linewidth}\raggedright
\(k\)
\end{minipage} & \begin{minipage}[b]{\linewidth}\raggedright
NLL
\end{minipage} & \begin{minipage}[b]{\linewidth}\raggedright
\(G^2\) (df, \(p\))
\end{minipage} & \begin{minipage}[b]{\linewidth}\raggedright
AIC
\end{minipage} & \begin{minipage}[b]{\linewidth}\raggedright
BIC
\end{minipage} & \begin{minipage}[b]{\linewidth}\raggedright
\(w_{\text{AIC}}\)
\end{minipage} & \begin{minipage}[b]{\linewidth}\raggedright
\(w_{\text{BIC}}\)
\end{minipage} & \begin{minipage}[b]{\linewidth}\raggedright
\(\hat d'_{1,2,4}\)
\end{minipage} & \begin{minipage}[b]{\linewidth}\raggedright
\(\hat\kappa\)
\end{minipage} \\
\midrule\noalign{}
\endfirsthead
\toprule\noalign{}
\begin{minipage}[b]{\linewidth}\raggedright
Model
\end{minipage} & \begin{minipage}[b]{\linewidth}\raggedright
\(k\)
\end{minipage} & \begin{minipage}[b]{\linewidth}\raggedright
NLL
\end{minipage} & \begin{minipage}[b]{\linewidth}\raggedright
\(G^2\) (df, \(p\))
\end{minipage} & \begin{minipage}[b]{\linewidth}\raggedright
AIC
\end{minipage} & \begin{minipage}[b]{\linewidth}\raggedright
BIC
\end{minipage} & \begin{minipage}[b]{\linewidth}\raggedright
\(w_{\text{AIC}}\)
\end{minipage} & \begin{minipage}[b]{\linewidth}\raggedright
\(w_{\text{BIC}}\)
\end{minipage} & \begin{minipage}[b]{\linewidth}\raggedright
\(\hat d'_{1,2,4}\)
\end{minipage} & \begin{minipage}[b]{\linewidth}\raggedright
\(\hat\kappa\)
\end{minipage} \\
\midrule\noalign{}
\endhead
\bottomrule\noalign{}
\endlastfoot
\(\mathcal{M}_0\) none & 2 & 27.27 & 35.55 (2, \textless{} .001) & 58.54
& 66.52 & 0.00 & 0.00 & 1.34, 1.34, 1.34 & 1.04 \\
\(\mathcal{M}_1\) log-linear & 3 & 9.77 & 0.54 (1, .46) & \textbf{25.53}
& \textbf{37.51} & \textbf{0.67} & \textbf{0.94} & 0.82, 1.37, 1.93 &
1.04 \\
\(\mathcal{M}_2\) free & 4 & 9.50 & 0 (saturated) & 26.99 & 42.96 & 0.33
& 0.06 & 0.78, 1.45, 1.88 & 1.04 \\
\end{longtable}
\end{landscape}

Likelihood-ratio tests: \(\mathcal{M}_0\) vs.~\(\mathcal{M}_1\):
\(\Lambda = 35.01\), df = 1, \(p < .001\); \(\mathcal{M}_1\)
vs.~\(\mathcal{M}_2\): \(\Lambda = 0.54\), df = 1, \(p = .46\).

\(r = 0.89\)

\textbf{Step 3: interpret.} Repetition clearly increases sensitivity:
the no-effect model misfits badly (\(G^2 = 35.5\), \(p < .001\)) and is
rejected by every criterion. The log-linear model fits as well as the
saturated model (LRT \(p = .46\)) with one parameter fewer, and both AIC
and BIC prefer it. Its estimate \(\hat\delta_1 \approx 0.55\) means each
doubling of study adds roughly half a \(d'\) unit. Recovery results show
that, at this design, both the parameters and the model choice can be
trusted. Caveats: this is one simulated participant; with only three
repetition levels, a log-linear trend cannot be distinguished from many
other increasing curves (e.g., a power function of \(r\)); and a single
criterion \(\kappa\) assumes the participant cannot tell how often an
item was studied before responding.

\begin{center}\rule{0.5\linewidth}{0.5pt}\end{center}

\Needspace{11\baselineskip}

\subsection{Lab (2 hours): midterm exam + project consultation}

\subsubsection{Midterm exam (90 minutes, in the lab slot)}

\begin{itemize}
\tightlist
\item
  \textbf{Coverage:} Weeks 1--7 (Chapters 1--3), with the emphasis on
  \emph{understanding and interpretation} (the course's specific goal
  2).
\item
  \textbf{Weight:} 3 of 20 points.
\item
  The midterm exam itself, with model answers, is in the course
  assessments (Stage 3 materials).
\end{itemize}

\textbf{Formula sheet provided at the exam:}

\[
\begin{aligned}
&d' = \Phi^{-1}(H) - \Phi^{-1}(F), \quad c = -\tfrac12[\Phi^{-1}(H) + \Phi^{-1}(F)], \quad H = \Phi(d'/2 - c), \quad F = \Phi(-d'/2 - c)\\
&f_{\exp}(t) = A e^{-\lambda t}, \quad f_{\text{pow}}(t) = A(1+t)^{-d}, \quad P_c^{\text{RW}} = \frac{1-r^z}{1-r^a},\ r = \frac{1-v}{1+v}, \quad \mathbb{E}[T] = \frac{aP_c - z}{v}\\
&\mathrm{SSE} = \textstyle\sum_i (y_i - \hat y_i)^2, \quad \mathrm{RMSD} = \sqrt{\mathrm{SSE}/N}\\
&\ell(\boldsymbol\theta) = \textstyle\sum_i \big[\ln\binom{n_i}{y_i} + y_i\ln f_i + (n_i - y_i)\ln(1-f_i)\big], \quad \mathrm{NLL} = -\ell\\
&G^2 = 2(\ell_{\text{sat}} - \hat\ell), \quad \Lambda = 2(\hat\ell_1 - \hat\ell_0) \sim \chi^2_{\Delta k}\\
&\mathrm{SE} \approx \sqrt{\mathrm{diag}(\mathbf{H}^{-1})}, \quad \text{Wald CI: } \hat\theta \pm 1.96\,\mathrm{SE}, \quad \text{profile CI: } 2(\mathrm{NLL}_{\text{prof}} - \mathrm{NLL}_{\min}) \le 3.84\\
&\mathrm{AIC} = 2\,\mathrm{NLL} + 2k, \quad \mathrm{AICc} = \mathrm{AIC} + \tfrac{2k(k+1)}{N-k-1}, \quad \mathrm{BIC} = 2\,\mathrm{NLL} + k\ln N, \quad w_i = \tfrac{e^{-\Delta_i/2}}{\sum_j e^{-\Delta_j/2}}
\end{aligned}
\]

Critical values: \(\chi^2_{1,.95} = 3.84\), \(\chi^2_{2,.95} = 5.99\),
\(\chi^2_{3,.95} = 7.81\). Quantiles: \(\Phi^{-1}(0.84) \approx 1.0\),
\(\Phi^{-1}(0.975) = 1.96\). Logs: \(\ln 2 = 0.693\),
\(\ln 10 = 2.303\), \(\ln 100 = 4.605\), \(\ln 400 = 5.991\).

\subsubsection{Project consultation (30 minutes, after the exam)}

Teams meet the instructor or TA for 5--10 minutes each to discuss their
project idea before the proposal deadline (start of Week 9). Bring a
one-paragraph description of your question, model and data.

\subsection{Connection to AI \& LLMs}

\textbf{Chatbots as study partners: trust, but verify.} Many students
use LLMs to explain concepts and check exam answers. That can help, but
LLMs can produce fluent, confident and wrong calculations. Treat a
chatbot's answer as you would a model's prediction: check it against
something you can compute. The chatbot used \(2\,\mathrm{NLL} + k\)
instead of \(2\,\mathrm{NLL} + 2k\).

The chatbot happens to pick the right winner (\(\mathcal{M}_2\)), but
every number is wrong, and it doubles the support for the most complex
model. A grader would give little credit, and in a real analysis the
error would favor overfitting.

What would BIC look like with the same error? - \ensuremath{\star}\ensuremath{\star} Write three small
``test cases'' (inputs with known answers) for the SDT formulas of Week
2 that you could use to check any answer, from a chatbot or a classmate.

\begin{center}\rule{0.5\linewidth}{0.5pt}\end{center}

\subsection{Common mistakes \& misconceptions (top 10 from Weeks 1--7)}

\begin{enumerate}
\def\labelenumi{\arabic{enumi}.}
\tightlist
\item
  Treating a good fit as evidence that a model is true (Weeks 3, 6).
\item
  Forgetting the minus sign: optimizers minimize the \textbf{negative}
  log-likelihood (Week 4).
\item
  Not clipping probabilities, so \(\ln 0\) crashes the fit (Week 4).
\item
  Least squares for counts or choices without thinking about the noise
  model (Weeks 3--4).
\item
  Reporting estimates without uncertainty, or using Wald intervals near
  bounds (Week 5).
\item
  Drawing conclusions about individuals from averaged curves (Week 5).
\item
  Leaking test data into training in cross-validation (Week 6).
\item
  Comparing non-nested models with an LRT, or reading AIC values on
  their own instead of as differences (Week 7).
\item
  Interpreting parameters or model comparisons without a recovery study
  at the actual design (Week 7).
\end{enumerate}

\begin{center}\rule{0.5\linewidth}{0.5pt}\end{center}

\subsection{Check your understanding (practice questions in exam style)}

\textbf{1. (Concept)} Explain in three sentences why the
exponential-with-asymptote model fitted the Week 2 data better than the
power model, although the power model generated the data.
\textbf{Answer} The asymptote model has an extra parameter, so it can
adapt to the particular noise of the sample, here the flat-looking
long-delay points. Better training fit is guaranteed for more flexible
models and is not evidence of truth. Generalization measures (CV,
AIC/BIC) and recovery studies are needed to judge it.

\textbf{2. (Calculation)} An observer has \(H = 0.84\) and \(F = 0.16\).
Compute \(d'\) and \(c\) (use \(\Phi^{-1}(0.84) \approx 1.0\)).
\textbf{Answer} \(d' = 1.0 - (-1.0) = 2.0\);
\(c = -\tfrac12(1.0 + (-1.0)) = 0\), so the observer is unbiased.

\textbf{3. (Calculation)} Two models are fitted to \(N = 100\) binary
responses. Model A: \(k = 2\), NLL = 60.0. Model B: \(k = 4\), NLL =
57.5. Compute AIC and BIC for both. Which does each criterion prefer?
(\(\ln 100 = 4.605\)) \textbf{Answer} AIC: A = 124.0, B = 123.0, so AIC
prefers B (slightly). BIC: A = 120 + 9.21 = 129.21; B = 115 + 18.42 =
133.42, so BIC prefers A. The criteria disagree because BIC penalizes
each parameter by 4.6 instead of 2.

\textbf{5. (Interpretation)} A recovery study gives
\(r(\text{true}, \text{recovered}) = 0.95\) for parameter \(A\) and
\(0.35\) for parameter \(B\). What do you conclude about a paper
reporting group differences in \(B\) at this design? \textbf{Answer}
\(B\) is poorly recoverable at this design: estimates are dominated by
noise or trade-offs, so group differences in estimated \(B\) are not
interpretable. The paper should improve the design or not interpret
\(B\).

\textbf{6. (Interpretation)} A BIC confusion matrix has P(selected =
power \textbar{} generated = asymptote) = 0.58. Your real data's BIC
winner is the power model. What should you write? \textbf{Answer} That
the power model is preferred by BIC but, at this design, BIC often
selects the power model for data generated by the asymptote model, so
the evidence does not discriminate between them well. Report the
inversion probabilities, AIC/CV results, and recommend a design with
longer delays.

\textbf{7. (Design)} You want to know whether a drug changes caution
(boundary \(a\)) or evidence quality (drift \(v\)) in the random walk
model. What data must you collect, and why is accuracy alone not enough?
\textbf{Answer} Choices \textbf{and} RTs (ideally whole RT
distributions) under drug and placebo. Both a larger \(a\) and a larger
\(v\) raise accuracy, but \(a\) slows responses while \(v\) speeds them,
so only the joint pattern separates them (Week 1).

\textbf{8. (Concept)} Why must cross-validation folds for the forgetting
data be formed within intervals if we want to test prediction of new
trials? \textbf{Answer} So that every training set contains all
retention intervals (conditions), and the test asks about new items in
known conditions. Leaving out whole intervals would test extrapolation
to new conditions, which is a different generalization claim.

\begin{center}\rule{0.5\linewidth}{0.5pt}\end{center}

\subsection{Assignment}

\begin{itemize}
\tightlist
\item
  \textbf{A3 (model comparison and recovery)} is due at the start of the
  Week 9 lecture.
\item
  \textbf{Project proposal} is due at the start of the Week 9 lecture
  (max. 2 pages; see Week 7).
\item
  No new assignment this week. Prepare for the midterm by running the
  case study above and answering the practice questions without looking
  at the answers.
\end{itemize}

\begin{center}\rule{0.5\linewidth}{0.5pt}\end{center}

\subsection{Further reading}

\begin{itemize}
\tightlist
\item
  \textbf{Farrell \& Lewandowsky (2018), chapters 2--5 and 10.} Reread
  the chapter summaries as a review of the whole first half of the
  course.
\item
  \textbf{Lee \& Wagenmakers (2014), chapter 11} (SDT). Read the model
  description as preparation for Week 9, where we start treating
  parameters as random variables with priors.
\end{itemize}

\emph{Preview for Week 9:} the second half of the course turns Bayesian.
Instead of a single best estimate plus a confidence interval, we will
compute the full probability distribution of the parameters given the
data.
